\Needspace{12\baselineskip}
\subsection{A non-Gauduchon background}
\label{sec:nongauduchon-weak-adjoints}
\index{Adjoint measure!non-Gauduchon background}

In the preceding example the background is closed, so the Gauduchon multipliers of all smooth positive residual forms are constant and the supporting probability is the normalized Monge--Amp\`ere measure of the residual. When no positive representative of the background is Gauduchon, the multiplier varies, and the supporting probability has to be identified through the adjoint of a linearized operator rather than through the mass. We now construct a divisorial reference in such a background on the same product surface. For continuous logarithmic densities depending only on the first factor, the solution has a continuous non-pluripolar residual. The weak Gauduchon equation for this residual has a unique normalized solution among nonnegative integrable weights, and the corresponding weighted measure is the only supporting probability, although the formula obtained for $B_\varphi$ covers only densities of this special form. The same weight is the limit of the classical multipliers along smooth approximations of the residual that keep the pole, whereas filling the pole produces mass on the divisor, as in \cref{thm:divisorial-smoothing-adjoints}.

We keep $X=\mathbb P^1\times\mathbb P^1$, the forms $\omega_1$, $\omega_2$ and $\omega=\omega_1+\omega_2$, the probability volume form $\mathrm dV=\omega_1\wedge\omega_2$, the affine coordinate $z$ on the first factor, $x=\log|z|^2$, $\chi(x)=\log(1+\mathrm e^x)$, $r=|z|^2/(1+|z|^2)$, $\psi=\log r$ and $D=\{0\}\times\mathbb P^1$ from the previous subsection, and we regard functions on either factor as functions on $X$. Fix $a,\lambda>0$, put $K=a+\lambda$, and let $b$ be a smooth, strictly positive, nonconstant function on the first factor. In this subsection the background and the residual background are the Hermitian forms
\[
 \theta=K\omega_1+b\omega_2,\qquad\beta=a\omega_1+b\omega_2.
\]
The Poincar\'e--Lelong formula $\ddc\psi=[D]-\omega_1$ gives
\begin{equation}\label{eq:nongauduchon-pole}
 \theta+\ddc(\lambda\psi)=\beta+\lambda[D].
\end{equation}
The forms of the previous subsection correspond to $K=2$, $\lambda=1$ and the constant function $b=2$.

On either factor, the Laplacian $\Delta$ is defined by $\ddc f=(\Delta f)\omega_i$ for smooth $f$. It is the Laplace--Beltrami operator of the round metric of the unit sphere, whose area form is $4\pi\omega_i$. For a smooth function $h$ on $X$, define $\Delta_1h$ and $\Delta_2h$ by
\[
 \ddc h\wedge\omega_2=(\Delta_1h)\,\mathrm dV,\qquad
 \ddc h\wedge\omega_1=(\Delta_2h)\,\mathrm dV;
\]
for a function on one factor, they are the Laplacian of that factor and zero. Since $K\omega_1$ and $\omega_2$ are closed,
\[
 \ddc\theta=\ddc\beta=\ddc b\wedge\omega_2=(\Delta b)\,\mathrm dV,
\]
and this form is not zero, because a function on $\mathbb P^1$ with $\Delta b=0$ is constant by the maximum principle. For every smooth real $v$, $\ddc(\theta+\ddc v)=\ddc\theta$ and $\ddc(\beta+\ddc v)=\ddc\beta$. Hence no smooth positive form $\theta+\ddc v$ or $\beta+\ddc v$ is Gauduchon, the Gauduchon multipliers of these forms are never constant, by the uniqueness in \cref{thm:gauduchon}, and \cref{prop:bounded-adjoint-persistence} applies to none of them.

\Needspace{10\baselineskip}
\begin{lemma}\label{lem:nongauduchon-divisorial-model}
Put $x_0=\log(\lambda/a)$ and $C_*=K\log K-a\log a-\lambda\log\lambda$. The envelope $\varphi=P_\theta[\lambda\psi]$ is
\begin{equation}\label{eq:nongauduchon-model}
 \varphi=W(x)-K\chi(x),\qquad
 W(x)=\begin{cases}K\chi(x_0)+\lambda(x-x_0),&x\leq x_0,\\K\chi(x),&x\geq x_0.\end{cases}
\end{equation}
It depends only on the first factor, satisfies $0\leq\varphi-\lambda\psi\leq C_*$, and satisfies the model condition \eqref{eq:als-model-condition} for every $\varepsilon$ with $\lambda/K<\varepsilon<1$.
\end{lemma}
\begin{proof}
Since $K\chi'(x_0)=K\lambda/(a+\lambda)=\lambda$, both one-sided derivatives of $W$ at $x_0$ equal $\lambda$, and the derivative of $K\chi$ increases on $[x_0,\infty)$; thus $W$ is convex. As $\omega_1=\ddc\chi(x)$ on the affine chart, $\varphi$ is $K\omega_1$-subharmonic there. It has a logarithmic pole with coefficient $\lambda$ at $z=0$ and vanishes where $|z|^2\geq\lambda/a$, so it is a $K\omega_1$-subharmonic function on the first factor, and $\varphi\in\PSH(X,\theta)$ because $b\omega_2\geq0$. Since $\lambda\psi=\lambda(x-\chi)$ and $\chi(x_0)=\log(K/a)$,
\[
 \varphi-\lambda\psi=
 \begin{cases}
 C_*-a\chi(x),&x\leq x_0,\\
 \lambda\log(1+\mathrm e^{-x}),&x\geq x_0.
 \end{cases}
\]
The first expression lies between $C_*-a\log(K/a)=\lambda\log(K/\lambda)$ and $C_*$, and the second between zero and $\lambda\log(K/\lambda)$, which proves $0\leq\varphi-\lambda\psi\leq C_*$ off $D$; on $D$ both functions equal $-\infty$.

Let $h\in\PSH(X,\theta)$ satisfy $h\leq0$ and $h\leq\lambda\psi+C$ for a constant $C$. Fix a point $y$ of the second factor. Since $\theta\leq K\omega_1+(\max_Xb)\,\omega_2$, which is closed, the sum of $h$, $K\log(1+|z|^2)$ and a local potential of $(\max_Xb)\,\omega_2$ near $y$ is plurisubharmonic near $\mathbb C\times\{y\}$. Its restriction to this line shows that $s_y(z)=h(z,y)+K\log(1+|z|^2)$ is subharmonic on $\mathbb C$ or identically $-\infty$. In the second case there is nothing to prove. In the first case, by the three-circles theorem, $M_y(x)=\max_{|z|^2=\mathrm e^x}s_y$ is a convex function of $x\in\mathbb R$, and the bounds on $h$ give
\[
 M_y(x)\leq K\chi(x),\qquad M_y(x)\leq\lambda x+a\chi(x)+C.
\]
As $x\to-\infty$, the second bound forces every one-sided slope of $M_y$ to be at least $\lambda$. Hence, for $x\leq x_0$,
\[
 M_y(x)\leq M_y(x_0)+\lambda(x-x_0)\leq K\chi(x_0)+\lambda(x-x_0)=W(x),
\]
while $M_y\leq K\chi=W$ for $x\geq x_0$. Thus $h\leq\varphi$ on the affine chart minus $D$. At $z=\infty$ we have $h\leq0=\varphi$, and on $D$ both are $-\infty$. Since $\varphi$ itself satisfies these bounds with $C=C_*$, the second description of $P_\theta[\cdot]$ in \cref{def:singularity-envelope} gives $P_\theta[\lambda\psi]=\varphi$. As $\varphi\sim\lambda\psi$ and the ordinary envelope depends only on the singularity type, also $P_\theta[\varphi]=\varphi$.

Finally, for $\lambda/K<\varepsilon<1$, \eqref{eq:nongauduchon-pole} gives
\[
 \varepsilon\theta+\ddc(\lambda\psi)
 =(\varepsilon K-\lambda)\omega_1+\varepsilon b\omega_2+\lambda[D]\geq0.
\]
Thus $\lambda\psi\in\PSH(X,\varepsilon\theta)$, and since $\lambda\psi\leq\varphi$, it is a competitor for $P_{\varepsilon\theta}(\varphi)$, which is therefore not identically $-\infty$.
\end{proof}

The lemma and its proof do not use that $b$ is nonconstant, and for $K=2$ and $\lambda=1$ the lemma is \cref{lem:divisorial-adjoint-model}. On the punctured disk $0<|z|^2<\lambda/a$ the function $W$ is affine in $x$, so that $\theta+\ddc\varphi=b\omega_2$ there: the residual of the reference degenerates on an open set. The solutions constructed below have strictly positive residual forms. As in the closed case, the ordinary type of $\varphi$ corresponds to the bounded potentials of the residual background.

\Needspace{10\baselineskip}
\begin{proposition}\label{prop:nongauduchon-residual-correspondence}
The map $w\mapsto\lambda\psi+w$ is a bijection from the bounded functions in $\PSH(X,\beta)$ onto the functions $v\in\PSH(X,\theta)$ with $v\sim\varphi$, and
\begin{equation}\label{eq:nongauduchon-residual-measure}
 \mathrm{MA}^{\mathrm{ext}}_\theta(\lambda\psi+w)=(\beta+\ddc w)^2.
\end{equation}
Consequently, $B_\varphi$ coincides on $L^\infty(\mathrm dV)$ with the normalization scalar of the Hermitian background $\beta$ and the zero reference, and for every $F\in C(X)$ the set $\mathcal A_\varphi(F)$ is the set of supporting probabilities of the latter at $F$. For every real $F\in L^\infty(\mathrm dV)$, the solution of \eqref{eq:relative-normalization} in the ordinary type of $\varphi$ is unique up to an additive constant.
\end{proposition}
\begin{proof}
Let $v\sim\varphi$. On $X\setminus D$, the function $w=v-\lambda\psi$ is bounded, by \cref{lem:nongauduchon-divisorial-model}, and $\beta+\ddc w=\theta+\ddc v\geq0$ there, by \eqref{eq:nongauduchon-pole}. By \cref{lem:qpsh-removable-pluripolar}, its upper-limit extension is a bounded function in $\PSH(X,\beta)$, and $v=\lambda\psi+w$ on $X$, since both sides equal $-\infty$ on $D$. Two such functions $w$ agree off $D$ and hence everywhere. Conversely, if $w\in\PSH(X,\beta)$ is bounded, then $\lambda\psi+w$ is $\theta$-plurisubharmonic, as the sum of a $\lambda\omega_1$-plurisubharmonic and a $\beta$-plurisubharmonic function, and it has the ordinary type of $\lambda\psi$, hence of $\varphi$.

Off $D$, the potential $v=\lambda\psi+w$ is locally bounded and the currents $\theta+\ddc v$ and $\beta+\ddc w$ agree, so their bounded products agree by \cref{prop:bounded-products}. Plurifine locality on the sets $\{v>-t\}\subset X\setminus D$ and the definition \eqref{eq:extended-ma} give $\mathrm{MA}^{\mathrm{ext}}_\theta(v)=\mathbf 1_{X\setminus D}(\beta+\ddc w)^2$, and the bounded product puts no mass on the pluripolar set $D$. This proves \eqref{eq:nongauduchon-residual-measure}.

The zero function is a model reference for the Hermitian form $\beta$. By the correspondence, a solution $v\sim\varphi$ of $\mathrm{MA}^{\mathrm{ext}}_\theta(v)=c\,\mathrm e^F\mathrm dV$ is the same as a bounded $\beta$-plurisubharmonic $w$ with $(\beta+\ddc w)^2=c\,\mathrm e^F\mathrm dV$, so the uniqueness of the constants identifies the two scalars, and hence their supporting sets. Finally, $\beta^2=2ab\,\mathrm dV$, so the density of $c\,\mathrm e^F\mathrm dV$ with respect to $\beta^2$ is bounded above and below by positive constants, and \cref{cor:bounded-positive-density-uniqueness} shows that two bounded solutions $w$ differ by a constant.
\end{proof}

The uniqueness of a supporting probability will follow from the uniqueness of a probability measure annihilating a linearized operator with merely continuous coefficients. Since the coefficients depend only on the first factor, an expansion in eigenfunctions along the second factor reduces this to a statement on the sphere.

\Needspace{10\baselineskip}
\begin{lemma}\label{lem:nongauduchon-invariant-probability}
Let $A$ be a continuous strictly positive function on the first factor, and let $\mu$ be a Borel probability measure on $X$ such that
\[
 \int_X\Bigl(\frac{\Delta_1h}{A}+\frac{\Delta_2h}{b}\Bigr)\mathrm d\mu=0
 \qquad\text{for every smooth real }h\text{ on }X.
\]
Then $\mu=A\,\mathrm dV/\int_{\mathbb P^1}A\,\omega_1$.
\end{lemma}
\begin{proof}
We first show that a finite signed Borel measure $T$ on $\mathbb P^1$ vanishes if $\int(\Delta f-Vf)\,\mathrm dT=0$ for every smooth $f$, where $V$ is continuous and strictly positive. Choose smooth functions $V_j\geq m$, with $m=\frac12\min V>0$, converging uniformly to $V$. Each operator $\Delta-V_j$ is elliptic and formally self-adjoint with respect to $\omega_1$, since $\int g\,\Delta f\,\omega_1=\int g\,\ddc f=\int f\,\ddc g$ by Stokes' theorem, and its kernel is zero by the maximum principle. By the Fredholm alternative and elliptic regularity, as in the proof of \cref{thm:gauduchon}, every smooth $g$ is of the form $\Delta f_j-V_jf_j$ with $f_j$ smooth. At a maximum point of $f_j$ we have $\Delta f_j\leq0$, hence $V_jf_j\leq-g$, and similarly at a minimum point; therefore $\|f_j\|_\infty\leq\|g\|_\infty/m$. Consequently
\[
 \Delta f_j-Vf_j=g+(V_j-V)f_j\longrightarrow g
\]
uniformly, and $\int g\,\mathrm dT=0$. Since smooth functions are uniformly dense in $C(\mathbb P^1)$, $T=0$.

Under the identification of $\mathbb P^1$ with the unit sphere, the restrictions of homogeneous harmonic polynomials of degree $\ell$ are eigenfunctions of $\Delta$ with eigenvalue $-\ell(\ell+1)$, and the sums of such restrictions are the restrictions of all polynomials, which are uniformly dense by the Stone--Weierstrass theorem. Choose real functions $e_0=1,e_1,e_2,\dots$ on the second factor, spanning this space, with $\Delta e_k=-\kappa_ke_k$ and $\kappa_k>0$ for $k\geq1$. Then $\int e_k\,\omega_2=-\kappa_k^{-1}\int\ddc e_k=0$ for $k\geq1$.

For each $k$, define a finite signed measure $T_k$ on the first factor by $\int f\,\mathrm dT_k=\int_Xfe_kA^{-1}\,\mathrm d\mu$. For smooth $f$ on the first factor, the function $h=fe_k$ satisfies $\Delta_1h=e_k\Delta f$ and $\Delta_2h=-\kappa_kfe_k$, so the hypothesis gives
\[
 \int\Bigl(\Delta f-\kappa_k\frac Abf\Bigr)\mathrm dT_k=0.
\]
For $k\geq1$, the first step with $V=\kappa_kA/b$ gives $T_k=0$. For $k=0$, the measure $T_0$ annihilates the Laplacians of smooth functions, which by the Fredholm alternative are exactly the smooth functions with zero $\omega_1$-integral; hence $T_0$ is a multiple of $\omega_1$, and testing with $f=A$ gives $T_0=\omega_1/\int A\,\omega_1$. Therefore, for every continuous $f$ on the first factor and every $k$,
\[
 \int_Xfe_k\,\mathrm d\mu=\int fA\,\mathrm dT_k
 =\frac{\int_XfA\,e_k\,\mathrm dV}{\int A\,\omega_1}.
\]
Finite sums of products of continuous functions on the two factors are uniformly dense in $C(X)$, and the functions on the second factor can be approximated uniformly by finite combinations of the $e_k$. Hence $\mu=A\,\mathrm dV/\int A\,\omega_1$.
\end{proof}

\Needspace{16\baselineskip}
\begin{theorem}\label{thm:nongauduchon-weak-adjoint}\index{Adjoint measure!weak linearization}
Let $F\in C(X)$ depend only on the first factor, and put
\[
 c=\frac{2a}{\int_Xb^{-1}\mathrm e^F\,\mathrm dV},\qquad
 A=\frac{c\,\mathrm e^F}{2b},\qquad
 \Omega=A\,\omega_1+b\,\omega_2.
\]
There is a continuous function $w$ on the first factor, unique up to an additive constant, with $\ddc w=(A-a)\omega_1$. The potential $u=\lambda\psi+w$ belongs to $\PSH(X,\theta)$, has the ordinary type of $\varphi$, and satisfies
\[
 \theta+\ddc u=\Omega+\lambda[D],\qquad
 \mathrm{MA}^{\mathrm{ext}}_\theta(u)=\Omega^2=c\,\mathrm e^F\,\mathrm dV.
\]
Every solution of \eqref{eq:relative-normalization} in the ordinary type of $\varphi$ is a translate of $u$, and
\[
 B_\varphi(F)=\log(2a)-\log\int_Xb^{-1}\mathrm e^F\,\mathrm dV.
\]
Moreover, the following hold.
\begin{enumerate}
\item For every smooth real $h$ and real $t$ with $\Omega+t\ddc h>0$, the potential $u+th$ has the ordinary type of $\varphi$, and $\mathrm{MA}^{\mathrm{ext}}_\theta(u+th)=(\Omega+t\ddc h)^2$. Its derivative at $t=0$ is $2\Omega\wedge\ddc h$, and the logarithmic linearization is
\[
 L_uh=\frac{2\Omega\wedge\ddc h}{\Omega^2}
 =\frac{\Delta_1h}{A}+\frac{\Delta_2h}{b}.
\]
\item The measure $\nu=\rho\Omega^2$, where $\rho=1/(2ab)$, is the only Borel probability measure on $X$ such that $\int_XL_uh\,\mathrm d\nu=0$ for every smooth real $h$, and
\[
 \mathcal A_\varphi(F)=\{\nu\},\qquad
 \nu=\frac{b^{-1}\mathrm e^F\,\mathrm dV}{\int_Xb^{-1}\mathrm e^F\,\mathrm dV}.
\]
It is not the normalized Monge--Amp\`ere measure of $u$.
\item The adjoint cofactor $Q=\rho\Omega$ satisfies $\mathrm dQ=0$ and $\int_X\Omega\wedge Q=1$. If a nonnegative measurable function $\eta$ satisfies $\int_X\eta\Omega^2=1$ and $\ddc(\eta\Omega)=0$ in the sense of currents, then $\eta=\rho$ almost everywhere.
\end{enumerate}
\end{theorem}
\begin{proof}
For $g$ on the first factor, $\int_Xg\,\mathrm dV=\int g\,\omega_1$, so $\int A\,\omega_1=a$. \Cref{thm:kn-continuous}, applied on the first factor, a curve with the Hermitian form $a\omega_1$, to the continuous density $A$, gives a constant $c'>0$ and a continuous $w$ with $a\omega_1+\ddc w=c'A\omega_1$. Integration over the first factor gives $c'=1$, and two such functions differ by a harmonic function on $\mathbb P^1$, hence by a constant. As a function on $X$, $w$ is bounded and $\beta$-plurisubharmonic, since $\beta+\ddc w=\Omega$ is positive. By \cref{prop:nongauduchon-residual-correspondence} and \eqref{eq:nongauduchon-pole}, $u\in\PSH(X,\theta)$, $u\sim\varphi$ and $\theta+\ddc u=\Omega+\lambda[D]$.

Let $h$ and $t$ be as in (1). Then $w+th$ is bounded and $\beta$-plurisubharmonic, and we compute its bounded product by writing $\beta+\ddc(w+th)=S+R_t$, where $S=a\omega_1+\ddc w=A\,\omega_1$ is a current with bounded potential pulled back from the first factor and $R_t=b\omega_2+t\ddc h$ is smooth. The products of \cref{prop:bounded-products} are multilinear; $S\wedge S=0$ by \cref{lem:bounded-products-submersion}, applied to the projection onto the first factor; and by the expansion \eqref{eq:bounded-expansion} with zero correction for the smooth factor, $S\wedge R_t$ is the product of a current of order zero by a smooth form. Hence
\[
 (\beta+\ddc(w+th))^2=2A\,\omega_1\wedge R_t+R_t^2=(\Omega+t\ddc h)^2,
\]
the last product being computed pointwise, since $\omega_1\wedge\omega_1=0$. By \cref{prop:nongauduchon-residual-correspondence}, $u+th\sim\varphi$ and $\mathrm{MA}^{\mathrm{ext}}_\theta(u+th)=(\Omega+t\ddc h)^2$. At $t=0$ this is $2Ab\,\mathrm dV=c\,\mathrm e^F\mathrm dV$, so $c_\varphi(F)=c$ by the uniqueness of the constant, which gives the formula for $B_\varphi(F)$; the solutions are translates of $u$ by \cref{prop:nongauduchon-residual-correspondence}. The expansion $(\Omega+t\ddc h)^2=\Omega^2+2t\,\Omega\wedge\ddc h+t^2(\ddc h)^2$ gives the derivative, and $\Omega\wedge\ddc h=(A\Delta_2h+b\Delta_1h)\,\mathrm dV$ gives the formula for $L_uh$.

The two terms of $Q=A(2ab)^{-1}\omega_1+(2a)^{-1}\omega_2$ are closed: the first is a continuous function on the first factor times a form of top degree on that factor, and $\omega_2$ is closed. Moreover $\nu=\rho\Omega^2=a^{-1}A\,\mathrm dV$, which is a probability measure equal to the expression in (2), and $\Omega\wedge Q=\nu$. Stokes' theorem for the closed current $Q$ gives
\[
 \int_XL_uh\,\mathrm d\nu=2\int_X\ddc h\wedge Q=0.
\]
By \cref{lem:nongauduchon-invariant-probability}, $\nu$ is the only probability measure with this property.

Let $\mu\in\mathcal A_\varphi(F)$ and let $h$ be smooth and real. For $|t|$ small, $\Omega+t\ddc h\geq\Omega/2$, and the function
\[
 F_t=F+\log\frac{(\Omega+t\ddc h)^2}{\Omega^2}
\]
is continuous. By (1), $u+th$ solves \eqref{eq:relative-normalization} at $F_t$ with the constant $c$, so $B_\varphi(F_t)=B_\varphi(F)$, and the supporting inequality \eqref{eq:adjoint-probability-set} at $G=F_t$ gives $\int_X(F_t-F)\,\mathrm d\mu\leq0$. Since $\Omega^2$ is a continuous positive multiple of $\mathrm dV$, we have $F_t-F=tL_uh+O(t^2)$ uniformly. Dividing by $t$ for both signs of $t$ and letting $t\to0$ shows that $\int_XL_uh\,\mathrm d\mu=0$, hence $\mu=\nu$. As $\mathcal A_\varphi(F)$ is nonempty by \cref{thm:weak-adjoint-probabilities}, it equals $\{\nu\}$. The normalized Monge--Amp\`ere measure of $u$ is $\mathrm e^F\mathrm dV/\int_X\mathrm e^F\mathrm dV$; it would agree with $\nu$ only if $b$ were constant.

Finally, let $\eta$ be as in (3). Since $\Omega^2=2Ab\,\mathrm dV$ and $2Ab$ is bounded below by a positive constant, $\eta$ is integrable, and $\eta\Omega$ is a current of order zero. For smooth real $h$,
\[
 \int_XL_uh\;\eta\Omega^2=2\int_X\eta\,\Omega\wedge\ddc h
 =2\int_Xh\,\ddc(\eta\Omega)=0.
\]
Thus $\eta\Omega^2$ is a probability measure with the property in (2), so $\eta\Omega^2=\rho\Omega^2$ and $\eta=\rho$ almost everywhere.
\end{proof}

Thus, in this example, the weak Gauduchon equation $\ddc(\rho\langle\theta+\ddc u\rangle)=0$ has a meaning: the non-pluripolar part $\Omega=\langle\theta+\ddc u\rangle$ of the current is a continuous form, its product with any integrable weight is defined, and the normalized nonnegative solution is unique. By \cref{prop:nongauduchon-residual-correspondence}, the theorem also describes the Hermitian background $\beta$ with the zero reference: for every continuous density depending only on the first factor, its supporting set consists of $\nu$ alone, and $\nu$ is the weighted measure of the weak adjoint of the continuous solution $w$. For another smooth positive volume form $\mathrm e^g\mathrm dV$, the scalar at $F$ is the scalar for $\mathrm dV$ at $F+g$, and the supporting sets correspond; the theorem therefore applies whenever $F+g$ depends only on the first factor.

\begin{example}\label{ex:nongauduchon-lipschitz-residual}
Let $a=1$ and $\lambda>0$, put $t=2r-1$, and take $b=2+t$ and $\mathrm e^F=2\bigl(\tfrac12+|t|\bigr)(2+t)$. The $\omega_1$-mass of $\{r<s\}$ is $s$ for $0\leq s\leq1$, so $t$ is uniformly distributed on $[-1,1]$ for $\mathrm dV$. Hence $c=1$, $A=\tfrac12+|t|$ and $\nu=\bigl(\tfrac12+|t|\bigr)\mathrm dV$. The residual coefficient $A$ and the density $F$ are Lipschitz but not differentiable along the circle $|z|=1$, where $t=0$ and $\mathrm dt\ne0$. Moreover $\int_Xt\,\mathrm d\nu=0$, whereas
\[
 \int_X\Omega^2=4,\qquad\int_Xt\,\frac{\Omega^2}4=\frac5{24}.
\]
\end{example}

The weak objects of \cref{thm:nongauduchon-weak-adjoint} are limits of the classical ones along smooth approximations of the residual that keep the pole along $D$. Since $\beta$ is not Gauduchon, the classical multipliers are not constant; the theorem needs neither bounds on their derivatives nor their pointwise convergence. For a real $(1,1)$-form $\alpha$, write $|\alpha|_\omega$ for its pointwise norm with respect to $\omega$.

\Needspace{14\baselineskip}
\begin{theorem}\label{thm:nongauduchon-fixed-pole-recovery}\index{Adjoint measure!fixed-divisor recovery}
Let $F$, $c$, $\Omega$, $\rho$, $Q$ and $\nu$ be as in \cref{thm:nongauduchon-weak-adjoint}. Let $w_j$ be smooth real functions on $X$ such that the forms $\Omega_j=\beta+\ddc w_j$ are positive and $\sup_X|\Omega_j-\Omega|_\omega\to0$. Let $\rho_j$ be the Gauduchon multiplier of $\Omega_j$ normalized by $\int_X\rho_j\Omega_j^2=1$, and put $Q_j=\rho_j\Omega_j$ and $\nu_j=\rho_j\Omega_j^2$. Then the potentials $u_j=\lambda\psi+w_j$ have the ordinary type of $\varphi$ and solve \eqref{eq:relative-normalization} with the constant $c$ at the smooth densities
\[
 F_j=\log\frac{\Omega_j^2}{c\,\mathrm dV},
\]
which converge uniformly to $F$, and $\mathcal A_\varphi(F_j)=\{\nu_j\}$. As $j\to\infty$,
\[
 \nu_j\rightharpoonup\nu,\qquad Q_j\rightharpoonup Q,\qquad
 \rho_j\,\mathrm dV\rightharpoonup\rho\,\mathrm dV.
\]
Such sequences exist, also with $w_j$ depending on both factors.
\end{theorem}
\begin{proof}
By \cref{prop:nongauduchon-residual-correspondence}, $u_j\sim\varphi$ and $\mathrm{MA}^{\mathrm{ext}}_\theta(u_j)=\Omega_j^2=c\,\mathrm e^{F_j}\mathrm dV$. The form $\Omega$ has continuous positive coefficients, so $\Omega\geq m\,\omega$ for some $m>0$, and for large $j$ also $\Omega_j\geq\frac m2\omega$. Hence the densities $q_j=\Omega_j^2/\mathrm dV$ converge uniformly to $q=\Omega^2/\mathrm dV$, all bounded below by a positive constant, and $F_j\to F$ uniformly. By \cref{prop:nongauduchon-residual-correspondence}, $\mathcal A_\varphi(F_j)$ is the supporting set at the smooth density $F_j$ for the Hermitian form $\beta$ and the zero reference; since $w_j$ is a smooth solution of that equation, \cref{thm:normalization-gauduchon-recovery} identifies it with $\{\nu_j\}$, the adjoint probability of $\Omega_j$ in \cref{prop:adjoint-gauduchon}. As $\mathcal A_\varphi(F)=\{\nu\}$ by \cref{thm:nongauduchon-weak-adjoint}, \cref{prop:weak-adjoint-stability} gives $\nu_j\rightharpoonup\nu$.

Let $\xi$ be a smooth real $(1,1)$-form; forms of type $(2,0)$ and $(0,2)$ pair trivially with the currents $Q_j$ and $Q$. The continuous functions $f_j=\Omega_j\wedge\xi/\Omega_j^2$ converge uniformly to $f=\Omega\wedge\xi/\Omega^2$, so
\[
 \int_XQ_j\wedge\xi=\int_Xf_j\,\mathrm d\nu_j\longrightarrow\int_Xf\,\mathrm d\nu=\int_XQ\wedge\xi.
\]
Similarly, for continuous $g$,
\[
 \int_Xg\rho_j\,\mathrm dV=\int_X\frac g{q_j}\,\mathrm d\nu_j
 \longrightarrow\int_X\frac gq\,\mathrm d\nu=\int_Xg\rho\,\mathrm dV.
\]

For the existence, choose smooth positive functions $A_j$ on the first factor with $\int A_j\omega_1=a$ converging uniformly to $A$, for instance normalized smooth uniform approximations of $A$. By the Fredholm alternative for the Laplacian there are smooth $w_j'$ on the first factor with $\ddc w_j'=(A_j-a)\omega_1$. For a smooth real function $k$ on $X$ and constants $\epsilon_j\to0$ so small that $A_j\omega_1+b\omega_2+\epsilon_j\ddc k>0$, the functions $w_j=w_j'+\epsilon_jk$ satisfy the hypotheses.
\end{proof}

The pole cannot be filled. As in \cref{thm:divisorial-smoothing-adjoints}, smooth potentials decreasing to $\lambda\psi$ keep a divisorial part in their residual forms, and in this background the limits of both their adjoint probabilities and their adjoint cofactors differ from the weak adjoint.

\Needspace{12\baselineskip}
\begin{proposition}\label{prop:nongauduchon-pole-filling}
For $0<\delta\leq1$, let $u_{\delta,0}=\log(r+\delta)$ as in \cref{thm:divisorial-smoothing-adjoints}, and put $\widehat\Omega_\delta=\theta+\ddc(\lambda u_{\delta,0})$. The smooth functions $\lambda u_{\delta,0}$ belong to $\PSH(X,\theta)$ and decrease to $\lambda\psi$, the forms $\widehat\Omega_\delta$ are positive, and their normalized Gauduchon multipliers are $\widehat\rho=1/(2Kb)$. As $\delta\downarrow0$,
\[
 \widehat\rho\,\widehat\Omega_\delta^2\rightharpoonup\frac aK\,\mathrm dV+\frac\lambda K[D]\wedge\omega_2,
 \qquad
 \widehat\rho\,\widehat\Omega_\delta\rightharpoonup\frac{a\,\omega_1}{2Kb}+\frac{\omega_2}{2K}+\frac{\lambda[D]}{2Kb(0)},
\]
where $b(0)$ is the value of $b$ at $z=0$. For $F_0=\log(2ab)$, \cref{thm:nongauduchon-weak-adjoint} gives $u=\lambda\psi$, $\mathcal A_\varphi(F_0)=\{\mathrm dV\}$ and $Q=\omega_1/(2b)+\omega_2/(2a)$; neither limit agrees with these.
\end{proposition}
\begin{proof}
On the affine chart, $r+\delta=((1+\delta)|z|^2+\delta)/(1+|z|^2)$, so the smooth closed form $\sigma_\delta=\omega_1+\ddc u_{\delta,0}$ on the first factor equals $\ddc\log(|z|^2+\epsilon)$ with $\epsilon=\delta/(1+\delta)$. It is positive, also near $z=\infty$, where it equals $\ddc\log(1+\epsilon|z|^{-2})$, and its mass on $\{|z|^2<s\}$ is $s/(s+\epsilon)$; hence $\int\sigma_\delta=1$ and $\sigma_\delta$ converges weakly to the Dirac measure at $z=0$. Therefore
\[
 \widehat\Omega_\delta=a\,\omega_1+\lambda\sigma_\delta+b\,\omega_2>0,
\]
so $\lambda u_{\delta,0}\in\PSH(X,\theta)$; these functions decrease to $\lambda\psi$ since $\log(r+\delta)$ decreases to $\log r$. The form $\widehat\rho\,\widehat\Omega_\delta=(a\omega_1+\lambda\sigma_\delta)/(2Kb)+\omega_2/(2K)$ is closed, since its first term is a smooth function on the first factor times a form of top degree on that factor. By \cref{thm:gauduchon}, $\widehat\rho$ is the Gauduchon multiplier of $\widehat\Omega_\delta$ up to a constant factor, and since $\widehat\Omega_\delta^2=2b\,(a\omega_1+\lambda\sigma_\delta)\wedge\omega_2$ and $\int(a\omega_1+\lambda\sigma_\delta)=K$, it is normalized. Thus $\widehat\rho\,\widehat\Omega_\delta^2=K^{-1}(a\omega_1+\lambda\sigma_\delta)\wedge\omega_2$, and the weak convergence of $\sigma_\delta$, together with the continuity of $1/b$, gives both limits. For $F_0$ we have $c=1$, $A=a$ and $w=0$, which gives the stated $u$, $\nu=\mathrm dV$ and $Q=\beta/(2ab)$. The first limit gives mass $\lambda/K$ to $D$, and the second has a divisorial part and a different coefficient of $\omega_2$.
\end{proof}

The forms $\widehat\Omega_\delta$ converge to $\beta+\lambda[D]$ only as currents, and the smooth potentials leave the ordinary type of $\varphi$; \cref{thm:nongauduchon-fixed-pole-recovery} does not apply to them. The proofs above use that $F$ depends only on the first factor, which makes the residual of the solution explicit. For a general continuous density, \cref{prop:nongauduchon-residual-correspondence} still gives a solution that is unique up to constants, but whether its supporting set is a singleton is the question of \cref{conj:continuous-gaud-adjoint} for the Hermitian form $\beta$.
